
\documentclass{article}
\usepackage{colortbl}
\usepackage{makecell}
\usepackage{multirow}
\usepackage{supertabular}

\begin{document}

\newcounter{utterance}

\centering \large Interaction Transcript for game `clean\_up', experiment `1\_medium\_3obj\_sl', episode 2 with deepseek{-}v4{-}pro.
\vspace{24pt}

{ \footnotesize  \setcounter{utterance}{1}
\setlength{\tabcolsep}{0pt}
\begin{supertabular}{c@{$\;$}|p{.15\linewidth}@{}p{.15\linewidth}p{.15\linewidth}p{.15\linewidth}p{.15\linewidth}p{.15\linewidth}}
    \# & $\;$A & \multicolumn{4}{c}{Game Master} & $\;\:$B\\
    \hline

    \theutterance \stepcounter{utterance}  
    & & \multicolumn{4}{p{0.6\linewidth}}{
        \cellcolor[rgb]{0.9,0.9,0.9}{
            \makecell[{{p{\linewidth}}}]{
                \texttt{\tiny{[P1$\langle$GM]}}
                \texttt{Jaz sem vaš vodja igre, vi pa igrate skupinsko igro z naslednjo mrežo kot igralno ploščo:} \\
\\ 
\texttt{```} \\
\texttt{ 1234567} \\
\texttt{╔═╤═╤═══╗} \\
\texttt{║◌│◌│C◌◌║ 1} \\
\texttt{╟─┴─┼─┬─╢ 2} \\
\texttt{║◌◌◌│◌│P║ 3} \\
\texttt{║◌◌◌├─┘◌║ 4} \\
\texttt{║◌┌─┘◌◌◌║ 5} \\
\texttt{╟─┤◌◌◌◌◌║ 6} \\
\texttt{║◌│◌L◌◌◌║ 7} \\
\texttt{╚═╧═════╝} \\
\\ 
\texttt{```} \\
\\ 
\texttt{* Na zgornjem robu so prikazane koordinate x, ki naraščajo proti desni, na desnem robu pa koordinate y, ki naraščajo navzdol.} \\
\texttt{* Na mrežo so naključno postavljeni naslednji predmeti: 'C', 'L', 'P'.} \\
\\ 
\\ 
\texttt{Drugi igralec vidi različico igralne plošče, na kateri so predmeti postavljeni na različna naključna mesta. Vi ne morete videti plošče drugega igralca, on pa ne more videti vaše.} \\
\\ 
\\ 
\texttt{**Cilj:**} \\
\\ 
\texttt{Oba igralca morata premikati predmete na svojem ozadju tako, da se enaki predmeti znajdejo na istih koordinatah. Z drugim igralcem se morate dogovoriti o skupni ciljni konfiguraciji.} \\
\\ 
\\ 
\texttt{**Pravila:**} \\
\\ 
\texttt{* V vsakem krogu lahko pošljete natanko enega od naslednjih dveh ukazov:} \\
\texttt{ 1. `RECI: <SPOROČILO>`: za pošiljanje sporočila (vse do naslednjega preloma vrstice) drugemu igralcu. Posredoval ga bom vašemu partnerju.} \\
\texttt{2. `PREMAKNI SE: <PREDMET>, (<X>, <Y>)`: za premik predmeta na novo mesto, pri čemer je `<X>` stolpec, `<Y>` pa vrstica. Obvestil vas bom, ali je bila vaša premestitev uspešna ali ne.} \\
\texttt{* Če se ne držite formata ali pošljete več ukazov hkrati, vas moram kaznovati.} \\
\texttt{* Če oba igralca zbereta več kot 8 kazni, oba izgubita igro.} \\
\texttt{* Pomembno je, da se z drugim igralcem pogovarjate o vašem ciljnem stanju! Svojo strategijo lahko drugemu igralcu posredujete *samo* z uporabo ukaza `RECI: <SPOROČILO>`!} \\
\\ 
\\ 
\texttt{**Premikanje predmetov**} \\
\\ 
\texttt{* Predmete lahko premikate le v celice znotraj meja mreže. Ciljna celica mora biti prazna, tj. vsebovati mora samo simbol '◌'.} \\
\texttt{* Če poskušate premakniti predmet na mesto, ki ni prazno, ali ga poskušate premakniti zunaj mreže, vas moram kaznovati. Dobili boste še en poskus.} \\
\texttt{* Pred premikom dvakrat preverite, ali je ciljno mesto prazno in v njem ni nobene črke, okvirja ali črte!} \\
\\ 
\\ 
\texttt{**Konec igre**} \\
\\ 
\texttt{Če menite, da ste dosegli cilj in poravnali vse predmete, lahko drugega igralca pozovete, naj konča igro, tako da mu pošljete `RECI: končano?`. Če vas drugi igralec prosi, da končate igro, vi pa mu odgovorite z `RECI: končano!`, se igra konča.} \\
\\ 
\texttt{Oba igralca zmagata, če se igra konča v 12 krogih, pri čemer je en krog definiran tako, da dva igralca pošljeta vsak po en veljaven ukaz.} \\
\\ 
\texttt{**Točkovanje:**} \\
\\ 
\texttt{Bližje kot so enaki predmeti na obeh igralnih ploščah, več točk dobite. Kazni zmanjšajo število točk. Ali lahko premagate rekord?} \\
\\ 
\\ 
\texttt{Pošljite sporočilo drugemu igralcu, da začnete igro!} \\
            }
        }
    }
    & \\ \\

    \theutterance \stepcounter{utterance}  
    & \multicolumn{4}{p{0.6\linewidth}}{
        \cellcolor[rgb]{0.9,0.9,0.9}{
            \makecell[{{p{\linewidth}}}]{
                \texttt{\tiny{[P1$\rangle$GM]}}
                \texttt{RECI: Pozdravljen! Začniva z usklajevanjem. Na moji plošči vidim predmete na teh lokacijah: C na (3,1), L na (3,7) in P na (7,3). Katere lokacije imaš ti?} \\
            }
        }
    }
    & & \\ \\

    \theutterance \stepcounter{utterance}  
    & & & \multicolumn{2}{p{0.3\linewidth}}{
        \cellcolor[rgb]{0.9,0.9,0.9}{
            \makecell[{{p{\linewidth}}}]{
                \texttt{\tiny{[GM$|$GM]}}
                \texttt{RECI: Pozdravljen! Začniva z usklajevanjem. Na moji plošči vidim predmete na teh lokacijah: C na (3,1), L na (3,7) in P na (7,3). Katere lokacije imaš ti?} \\
            }
        }
    }
    & & \\ \\

    \theutterance \stepcounter{utterance}  
    & & \multicolumn{4}{p{0.6\linewidth}}{
        \cellcolor[rgb]{0.9,0.9,0.9}{
            \makecell[{{p{\linewidth}}}]{
                \texttt{\tiny{[GM$\rangle$P2]}}
                \texttt{Jaz sem vaš vodja igre, vi pa igrate skupinsko igro z naslednjo mrežo kot igralno ploščo:} \\
\\ 
\texttt{```} \\
\texttt{ 1234567} \\
\texttt{╔═╤═╤═══╗} \\
\texttt{║◌│◌│◌◌◌║ 1} \\
\texttt{╟─┴─┼─┬─╢ 2} \\
\texttt{║◌◌◌│◌│L║ 3} \\
\texttt{║C◌◌├─┘P║ 4} \\
\texttt{║◌┌─┘◌◌◌║ 5} \\
\texttt{╟─┤◌◌◌◌◌║ 6} \\
\texttt{║◌│◌◌◌◌◌║ 7} \\
\texttt{╚═╧═════╝} \\
\\ 
\texttt{```} \\
\\ 
\texttt{* Na zgornjem robu so prikazane koordinate x, ki naraščajo proti desni, na desnem robu pa koordinate y, ki naraščajo navzdol.} \\
\texttt{* Na mrežo so naključno postavljeni naslednji predmeti: 'C', 'L', 'P'.} \\
\\ 
\\ 
\texttt{Drugi igralec vidi različico igralne plošče, na kateri so predmeti postavljeni na različna naključna mesta. Vi ne morete videti plošče drugega igralca, on pa ne more videti vaše.} \\
\\ 
\\ 
\texttt{**Cilj:**} \\
\\ 
\texttt{Oba igralca morata premikati predmete na svojem ozadju tako, da se enaki predmeti znajdejo na istih koordinatah. Z drugim igralcem se morate dogovoriti o skupni ciljni konfiguraciji.} \\
\\ 
\\ 
\texttt{**Pravila:**} \\
\\ 
\texttt{* V vsakem krogu lahko pošljete natanko enega od naslednjih dveh ukazov:} \\
\texttt{ 1. `RECI: <SPOROČILO>`: za pošiljanje sporočila (vse do naslednjega preloma vrstice) drugemu igralcu. Posredoval ga bom vašemu partnerju.} \\
\texttt{2. `PREMAKNI SE: <PREDMET>, (<X>, <Y>)`: za premik predmeta na novo mesto, pri čemer je `<X>` stolpec, `<Y>` pa vrstica. Obvestil vas bom, ali je bila vaša premestitev uspešna ali ne.} \\
\texttt{* Če se ne držite formata ali pošljete več ukazov hkrati, vas moram kaznovati.} \\
\texttt{* Če oba igralca zbereta več kot 8 kazni, oba izgubita igro.} \\
\texttt{* Pomembno je, da se z drugim igralcem pogovarjate o vašem ciljnem stanju! Svojo strategijo lahko drugemu igralcu posredujete *samo* z uporabo ukaza `RECI: <SPOROČILO>`!} \\
\\ 
\\ 
\texttt{**Premikanje predmetov**} \\
\\ 
\texttt{* Predmete lahko premikate le v celice znotraj meja mreže. Ciljna celica mora biti prazna, tj. vsebovati mora samo simbol '◌'.} \\
\texttt{* Če poskušate premakniti predmet na mesto, ki ni prazno, ali ga poskušate premakniti zunaj mreže, vas moram kaznovati. Dobili boste še en poskus.} \\
\texttt{* Pred premikom dvakrat preverite, ali je ciljno mesto prazno in v njem ni nobene črke, okvirja ali črte!} \\
\\ 
\\ 
\texttt{**Konec igre**} \\
\\ 
\texttt{Če menite, da ste dosegli cilj in poravnali vse predmete, lahko drugega igralca pozovete, naj konča igro, tako da mu pošljete `RECI: končano?`. Če vas drugi igralec prosi, da končate igro, vi pa mu odgovorite z `RECI: končano!`, se igra konča.} \\
\\ 
\texttt{Oba igralca zmagata, če se igra konča v 12 krogih, pri čemer je en krog definiran tako, da dva igralca pošljeta vsak po en veljaven ukaz.} \\
\\ 
\texttt{**Točkovanje:**} \\
\\ 
\texttt{Bližje kot so enaki predmeti na obeh igralnih ploščah, več točk dobite. Kazni zmanjšajo število točk. Ali lahko premagate rekord?} \\
\\ 
\texttt{Drugi igralec je začel igro s tem sporočilom:} \\
\\ 
\texttt{"Pozdravljen! Začniva z usklajevanjem. Na moji plošči vidim predmete na teh lokacijah: C na (3,1), L na (3,7) in P na (7,3). Katere lokacije imaš ti?"} \\
\\ 
\texttt{Kakšen je vaš prvi ukaz?} \\
            }
        }
    }
    & \\ \\

    \theutterance \stepcounter{utterance}  
    & & & \multicolumn{4}{p{0.6\linewidth}}{
        \cellcolor[rgb]{0.9,0.9,0.9}{
            \makecell[{{p{\linewidth}}}]{
                \texttt{\tiny{[GM$\langle$P2]}}
                \texttt{RECI: Pozdravljen! Na moji plošči so predmeti: C na (2,4), L na (7,3) in P na (7,4). Predlagam, da se dogovoriva za skupne ciljne lokacije. Kaj meniš o tem, da damo C na (3,3), L na (4,4) in P na (5,5)? Ali imaš kakšno boljšo idejo?} \\
            }
        }
    }
     \\ \\

    \theutterance \stepcounter{utterance}  
    & & & \multicolumn{2}{p{0.3\linewidth}}{
        \cellcolor[rgb]{0.9,0.9,0.9}{
            \makecell[{{p{\linewidth}}}]{
                \texttt{\tiny{[GM$|$GM]}}
                \texttt{RECI: Pozdravljen! Na moji plošči so predmeti: C na (2,4), L na (7,3) in P na (7,4). Predlagam, da se dogovoriva za skupne ciljne lokacije. Kaj meniš o tem, da damo C na (3,3), L na (4,4) in P na (5,5)? Ali imaš kakšno boljšo idejo?} \\
            }
        }
    }
    & & \\ \\

    \theutterance \stepcounter{utterance}  
    & & \multicolumn{4}{p{0.6\linewidth}}{
        \cellcolor[rgb]{0.9,0.9,0.9}{
            \makecell[{{p{\linewidth}}}]{
                \texttt{\tiny{[P1$\langle$GM]}}
                \texttt{Vaše sporočilo je bilo posredovano drugemu igralcu.} \\
\\ 
\\ 
\texttt{Trenutno igrate krog 1 od največ 12.} \\
\texttt{Skupaj ste zbrali 0 od 8 kazni.} \\
\texttt{Drugi igralec je poslal to sporočilo:} \\
\\ 
\\ 
\texttt{"Pozdravljen! Na moji plošči so predmeti: C na (2,4), L na (7,3) in P na (7,4). Predlagam, da se dogovoriva za skupne ciljne lokacije. Kaj meniš o tem, da damo C na (3,3), L na (4,4) in P na (5,5)? Ali imaš kakšno boljšo idejo?"} \\
\\ 
\\ 
\texttt{Kakšen je vaš naslednji ukaz?} \\
            }
        }
    }
    & \\ \\

    \theutterance \stepcounter{utterance}  
    & \multicolumn{4}{p{0.6\linewidth}}{
        \cellcolor[rgb]{0.9,0.9,0.9}{
            \makecell[{{p{\linewidth}}}]{
                \texttt{\tiny{[P1$\rangle$GM]}}
                \texttt{RECI: Dober predlog, ampak preveriva, če so te lokacije proste na obeh ploščah. Na moji plošči: C je trenutno na (3,1), L na (3,7), P na (7,3). Tvoje ciljne lokacije (3,3), (4,4) in (5,5) so pri meni prazne, tako da jih lahko dosežem. Preden potrdiva, preveri še ti, če so (3,3), (4,4) in (5,5) prazne na tvoji plošči. Če so, se strinjam s tem ciljem!} \\
            }
        }
    }
    & & \\ \\

    \theutterance \stepcounter{utterance}  
    & & & \multicolumn{2}{p{0.3\linewidth}}{
        \cellcolor[rgb]{0.9,0.9,0.9}{
            \makecell[{{p{\linewidth}}}]{
                \texttt{\tiny{[GM$|$GM]}}
                \texttt{RECI: Dober predlog, ampak preveriva, če so te lokacije proste na obeh ploščah. Na moji plošči: C je trenutno na (3,1), L na (3,7), P na (7,3). Tvoje ciljne lokacije (3,3), (4,4) in (5,5) so pri meni prazne, tako da jih lahko dosežem. Preden potrdiva, preveri še ti, če so (3,3), (4,4) in (5,5) prazne na tvoji plošči. Če so, se strinjam s tem ciljem!} \\
            }
        }
    }
    & & \\ \\

    \theutterance \stepcounter{utterance}  
    & & \multicolumn{4}{p{0.6\linewidth}}{
        \cellcolor[rgb]{0.9,0.9,0.9}{
            \makecell[{{p{\linewidth}}}]{
                \texttt{\tiny{[GM$\rangle$P2]}}
                \texttt{Vaše sporočilo je bilo posredovano drugemu igralcu.} \\
\\ 
\\ 
\texttt{Trenutno igrate krog 2 od največ 12.} \\
\texttt{Skupaj ste zbrali 0 od 8 kazni.} \\
\texttt{Drugi igralec je poslal to sporočilo:} \\
\\ 
\\ 
\texttt{"Dober predlog, ampak preveriva, če so te lokacije proste na obeh ploščah. Na moji plošči: C je trenutno na (3,1), L na (3,7), P na (7,3). Tvoje ciljne lokacije (3,3), (4,4) in (5,5) so pri meni prazne, tako da jih lahko dosežem. Preden potrdiva, preveri še ti, če so (3,3), (4,4) in (5,5) prazne na tvoji plošči. Če so, se strinjam s tem ciljem!"} \\
\\ 
\\ 
\texttt{Kakšen je vaš naslednji ukaz?} \\
            }
        }
    }
    & \\ \\

    \theutterance \stepcounter{utterance}  
    & & & \multicolumn{4}{p{0.6\linewidth}}{
        \cellcolor[rgb]{0.9,0.9,0.9}{
            \makecell[{{p{\linewidth}}}]{
                \texttt{\tiny{[GM$\langle$P2]}}
                \texttt{Najprej preverim ciljne lokacije na svoji plošči:} \\
\\ 
\texttt{{-} (3,3): v vrstici 3, stolpec 3 je '◌' (prazno) ✓} \\
\texttt{{-} (4,4): v vrstici 4, stolpec 4 je '◌' (prazno) ✓} \\
\texttt{{-} (5,5): v vrstici 5, stolpec 5 je '◌' (prazno) ✓} \\
\\ 
\texttt{Vse tri ciljne lokacije so na moji plošči prazne. Zdaj pa premaknem prvi predmet. Začel bom s C, ki je trenutno na (2,4), proti cilju na (3,3). Ampak najprej se moram prepričati, da je pot prosta {-} ker premikamo direktno na novo lokacijo, ne rabim skrbeti za vmesne korake.} \\
\\ 
\texttt{PREMAKNI SE: C, (3,3)} \\
            }
        }
    }
     \\ \\

    \theutterance \stepcounter{utterance}  
    & & & \multicolumn{2}{p{0.3\linewidth}}{
        \cellcolor[rgb]{0.9,0.9,0.9}{
            \makecell[{{p{\linewidth}}}]{
                \texttt{\tiny{[GM$|$GM]}}
                \texttt{Najprej preverim ciljne lokacije na svoji plošči:} \\
\\ 
\texttt{{-} (3,3): v vrstici 3, stolpec 3 je '◌' (prazno) ✓} \\
\texttt{{-} (4,4): v vrstici 4, stolpec 4 je '◌' (prazno) ✓} \\
\texttt{{-} (5,5): v vrstici 5, stolpec 5 je '◌' (prazno) ✓} \\
\\ 
\texttt{Vse tri ciljne lokacije so na moji plošči prazne. Zdaj pa premaknem prvi predmet. Začel bom s C, ki je trenutno na (2,4), proti cilju na (3,3). Ampak najprej se moram prepričati, da je pot prosta {-} ker premikamo direktno na novo lokacijo, ne rabim skrbeti za vmesne korake.} \\
\\ 
\texttt{PREMAKNI SE: C, (3,3)} \\
            }
        }
    }
    & & \\ \\

    \theutterance \stepcounter{utterance}  
    & & & \multicolumn{2}{p{0.3\linewidth}}{
        \cellcolor[rgb]{0.9,0.9,0.9}{
            \makecell[{{p{\linewidth}}}]{
                \texttt{\tiny{[GM$|$GM]}}
                \texttt{Invalid format: head is not empty: 'Najprej preverim ciljne lokacije na svoji plošči:} \\
\\ 
\texttt{{-} (3,3): v vrstici 3, stolpec 3 je '◌' (prazno) ✓} \\
\texttt{{-} (4,4): v vrstici 4, stolpec 4 je '◌' (prazno) ✓} \\
\texttt{{-} (5,5): v vrstici 5, stolpec 5 je '◌' (prazno) ✓} \\
\\ 
\texttt{Vse tri ciljne lokacije so na moji plošči prazne. Zdaj pa premaknem prvi predmet. Začel bom s C, ki je trenutno na (2,4), proti cilju na (3,3). Ampak najprej se moram prepričati, da je pot prosta {-} ker premikamo direktno na novo lokacijo, ne rabim skrbeti za vmesne korake.} \\
\\ 
\texttt{'} \\
            }
        }
    }
    & & \\ \\

    \theutterance \stepcounter{utterance}  
    & & \multicolumn{4}{p{0.6\linewidth}}{
        \cellcolor[rgb]{0.9,0.9,0.9}{
            \makecell[{{p{\linewidth}}}]{
                \texttt{\tiny{[GM$\rangle$P2]}}
                \texttt{Kazen: Vaše sporočilo ne sme vsebovati ničesar pred ukazom!} \\
\\ 
\texttt{Prepričajte se, da vaš odgovor vsebuje samo `RECI: <SPOROČILO>` ali `PREMAKNI SE: <PREDMET>, (<X>, <Y>)` in nič drugega!} \\
\\ 
\texttt{Skupaj ste zbrali 1 od 8 kazni.} \\
\texttt{Prosim, poskusite znova!} \\
            }
        }
    }
    & \\ \\

    \theutterance \stepcounter{utterance}  
    & & & \multicolumn{4}{p{0.6\linewidth}}{
        \cellcolor[rgb]{0.9,0.9,0.9}{
            \makecell[{{p{\linewidth}}}]{
                \texttt{\tiny{[GM$\langle$P2]}}
                \texttt{PREMAKNI SE: C, (3,3)} \\
            }
        }
    }
     \\ \\

    \theutterance \stepcounter{utterance}  
    & & & \multicolumn{2}{p{0.3\linewidth}}{
        \cellcolor[rgb]{0.9,0.9,0.9}{
            \makecell[{{p{\linewidth}}}]{
                \texttt{\tiny{[GM$|$GM]}}
                \texttt{PREMAKNI SE: C, (3,3)} \\
            }
        }
    }
    & & \\ \\

    \theutterance \stepcounter{utterance}  
    & & & \multicolumn{2}{p{0.3\linewidth}}{
        \cellcolor[rgb]{0.9,0.9,0.9}{
            \makecell[{{p{\linewidth}}}]{
                \texttt{\tiny{[GM$|$GM]}}
                \texttt{C moved to (3, 3)} \\
            }
        }
    }
    & & \\ \\

    \theutterance \stepcounter{utterance}  
    & & \multicolumn{4}{p{0.6\linewidth}}{
        \cellcolor[rgb]{0.9,0.9,0.9}{
            \makecell[{{p{\linewidth}}}]{
                \texttt{\tiny{[P1$\langle$GM]}}
                \texttt{Vaše sporočilo je bilo posredovano drugemu igralcu.} \\
\\ 
\\ 
\texttt{Trenutno igrate krog 2 od največ 12.} \\
\texttt{Skupaj ste zbrali 1 od 8 kazni.} \\
\texttt{Drugi igralec je premaknil predmet na svoji mreži.} \\
\\ 
\\ 
\texttt{Kakšen je vaš naslednji ukaz?} \\
            }
        }
    }
    & \\ \\

    \theutterance \stepcounter{utterance}  
    & \multicolumn{4}{p{0.6\linewidth}}{
        \cellcolor[rgb]{0.9,0.9,0.9}{
            \makecell[{{p{\linewidth}}}]{
                \texttt{\tiny{[P1$\rangle$GM]}}
                \texttt{PREMAKNI SE: C, (3,2)} \\
            }
        }
    }
    & & \\ \\

    \theutterance \stepcounter{utterance}  
    & & & \multicolumn{2}{p{0.3\linewidth}}{
        \cellcolor[rgb]{0.9,0.9,0.9}{
            \makecell[{{p{\linewidth}}}]{
                \texttt{\tiny{[GM$|$GM]}}
                \texttt{PREMAKNI SE: C, (3,2)} \\
            }
        }
    }
    & & \\ \\

    \theutterance \stepcounter{utterance}  
    & & & \multicolumn{2}{p{0.3\linewidth}}{
        \cellcolor[rgb]{0.9,0.9,0.9}{
            \makecell[{{p{\linewidth}}}]{
                \texttt{\tiny{[GM$|$GM]}}
                \texttt{Kazen: (3,2) ni prazen, vendar vsebuje '─'.} \\
\\ 
\texttt{Skupaj ste zbrali 2 od 8 kazni.} \\
\texttt{Prosim, poskusite znova!} \\
            }
        }
    }
    & & \\ \\

    \theutterance \stepcounter{utterance}  
    & & \multicolumn{4}{p{0.6\linewidth}}{
        \cellcolor[rgb]{0.9,0.9,0.9}{
            \makecell[{{p{\linewidth}}}]{
                \texttt{\tiny{[P1$\langle$GM]}}
                \texttt{Kazen: (3,2) ni prazen, vendar vsebuje '─'.} \\
\\ 
\texttt{Skupaj ste zbrali 2 od 8 kazni.} \\
\texttt{Prosim, poskusite znova!} \\
            }
        }
    }
    & \\ \\

    \theutterance \stepcounter{utterance}  
    & \multicolumn{4}{p{0.6\linewidth}}{
        \cellcolor[rgb]{0.9,0.9,0.9}{
            \makecell[{{p{\linewidth}}}]{
                \texttt{\tiny{[P1$\rangle$GM]}}
                \texttt{PREMAKNI SE: C, (2,1)} \\
            }
        }
    }
    & & \\ \\

    \theutterance \stepcounter{utterance}  
    & & & \multicolumn{2}{p{0.3\linewidth}}{
        \cellcolor[rgb]{0.9,0.9,0.9}{
            \makecell[{{p{\linewidth}}}]{
                \texttt{\tiny{[GM$|$GM]}}
                \texttt{PREMAKNI SE: C, (2,1)} \\
            }
        }
    }
    & & \\ \\

    \theutterance \stepcounter{utterance}  
    & & & \multicolumn{2}{p{0.3\linewidth}}{
        \cellcolor[rgb]{0.9,0.9,0.9}{
            \makecell[{{p{\linewidth}}}]{
                \texttt{\tiny{[GM$|$GM]}}
                \texttt{Kazen: (2,1) ni prazen, vendar vsebuje '│'.} \\
\\ 
\texttt{Skupaj ste zbrali 3 od 8 kazni.} \\
\texttt{Prosim, poskusite znova!} \\
            }
        }
    }
    & & \\ \\

    \theutterance \stepcounter{utterance}  
    & & \multicolumn{4}{p{0.6\linewidth}}{
        \cellcolor[rgb]{0.9,0.9,0.9}{
            \makecell[{{p{\linewidth}}}]{
                \texttt{\tiny{[P1$\langle$GM]}}
                \texttt{Kazen: (2,1) ni prazen, vendar vsebuje '│'.} \\
\\ 
\texttt{Skupaj ste zbrali 3 od 8 kazni.} \\
\texttt{Prosim, poskusite znova!} \\
            }
        }
    }
    & \\ \\

    \theutterance \stepcounter{utterance}  
    & \multicolumn{4}{p{0.6\linewidth}}{
        \cellcolor[rgb]{0.9,0.9,0.9}{
            \makecell[{{p{\linewidth}}}]{
                \texttt{\tiny{[P1$\rangle$GM]}}
                \texttt{PREMAKNI SE: C, (1,1)} \\
            }
        }
    }
    & & \\ \\

    \theutterance \stepcounter{utterance}  
    & & & \multicolumn{2}{p{0.3\linewidth}}{
        \cellcolor[rgb]{0.9,0.9,0.9}{
            \makecell[{{p{\linewidth}}}]{
                \texttt{\tiny{[GM$|$GM]}}
                \texttt{PREMAKNI SE: C, (1,1)} \\
            }
        }
    }
    & & \\ \\

    \theutterance \stepcounter{utterance}  
    & & & \multicolumn{2}{p{0.3\linewidth}}{
        \cellcolor[rgb]{0.9,0.9,0.9}{
            \makecell[{{p{\linewidth}}}]{
                \texttt{\tiny{[GM$|$GM]}}
                \texttt{C moved to (1, 1)} \\
            }
        }
    }
    & & \\ \\

    \theutterance \stepcounter{utterance}  
    & & \multicolumn{4}{p{0.6\linewidth}}{
        \cellcolor[rgb]{0.9,0.9,0.9}{
            \makecell[{{p{\linewidth}}}]{
                \texttt{\tiny{[GM$\rangle$P2]}}
                \texttt{Uspešno premaknjen 'C' v (3,3). Vaša posodobljena mreža je videti takole:} \\
\texttt{```} \\
\texttt{ 1234567} \\
\texttt{╔═╤═╤═══╗} \\
\texttt{║◌│◌│◌◌◌║ 1} \\
\texttt{╟─┴─┼─┬─╢ 2} \\
\texttt{║◌◌C│◌│L║ 3} \\
\texttt{║◌◌◌├─┘P║ 4} \\
\texttt{║◌┌─┘◌◌◌║ 5} \\
\texttt{╟─┤◌◌◌◌◌║ 6} \\
\texttt{║◌│◌◌◌◌◌║ 7} \\
\texttt{╚═╧═════╝} \\
\texttt{```} \\
\texttt{Trenutno igrate krog 3 od največ 12.} \\
\texttt{Skupaj ste zbrali 3 od 8 kazni.} \\
\texttt{Drugi igralec je premaknil predmet na svoji mreži.} \\
\\ 
\\ 
\texttt{Kakšen je vaš naslednji ukaz?} \\
            }
        }
    }
    & \\ \\

    \theutterance \stepcounter{utterance}  
    & & & \multicolumn{4}{p{0.6\linewidth}}{
        \cellcolor[rgb]{0.9,0.9,0.9}{
            \makecell[{{p{\linewidth}}}]{
                \texttt{\tiny{[GM$\langle$P2]}}
                \texttt{PREMAKNI SE: L, (4,4)} \\
            }
        }
    }
     \\ \\

    \theutterance \stepcounter{utterance}  
    & & & \multicolumn{2}{p{0.3\linewidth}}{
        \cellcolor[rgb]{0.9,0.9,0.9}{
            \makecell[{{p{\linewidth}}}]{
                \texttt{\tiny{[GM$|$GM]}}
                \texttt{PREMAKNI SE: L, (4,4)} \\
            }
        }
    }
    & & \\ \\

    \theutterance \stepcounter{utterance}  
    & & & \multicolumn{2}{p{0.3\linewidth}}{
        \cellcolor[rgb]{0.9,0.9,0.9}{
            \makecell[{{p{\linewidth}}}]{
                \texttt{\tiny{[GM$|$GM]}}
                \texttt{Kazen: (4,4) ni prazen, vendar vsebuje '├'.} \\
\\ 
\texttt{Skupaj ste zbrali 4 od 8 kazni.} \\
\texttt{Prosim, poskusite znova!} \\
            }
        }
    }
    & & \\ \\

    \theutterance \stepcounter{utterance}  
    & & \multicolumn{4}{p{0.6\linewidth}}{
        \cellcolor[rgb]{0.9,0.9,0.9}{
            \makecell[{{p{\linewidth}}}]{
                \texttt{\tiny{[GM$\rangle$P2]}}
                \texttt{Kazen: (4,4) ni prazen, vendar vsebuje '├'.} \\
\\ 
\texttt{Skupaj ste zbrali 4 od 8 kazni.} \\
\texttt{Prosim, poskusite znova!} \\
            }
        }
    }
    & \\ \\

    \theutterance \stepcounter{utterance}  
    & & & \multicolumn{4}{p{0.6\linewidth}}{
        \cellcolor[rgb]{0.9,0.9,0.9}{
            \makecell[{{p{\linewidth}}}]{
                \texttt{\tiny{[GM$\langle$P2]}}
                \texttt{Ponovno preverim ploščo. L je trenutno na (7,3), cilj je (4,4). Na (4,4) je res '├', kar ni prazno. Ciljne lokacije smo določili kot (3,3) za C, (4,4) za L in (5,5) za P. Ampak (4,4) ni prazno. Počakaj, ponovno preverim svojo trenutno ploščo.} \\
\\ 
\texttt{Vrstica 4: `║◌◌◌├─┘P║ 4` {-} stolpci: 1:◌, 2:◌, 3:◌, 4:├, 5:─, 6:┘, 7:P. Torej (4,4) je res zasedeno z '├'. To pomeni, da (4,4) ni dosegljivo. } \\
\\ 
\texttt{Drugi igralec je rekel, da so pri njem te lokacije prazne, ampak pri meni (4,4) ni prazno. Moram predlagati nove ciljne lokacije.} \\
\\ 
\texttt{RECI: Oprosti, na moji plošči (4,4) ni prazno {-} tam je '├'. Prav tako (5,5) preverim... Vrstica 5: `║◌┌─┘◌◌◌║ 5` {-} (5,5) je ◌ (prazno). (3,3) je zdaj C, kar je v redu. Predlagam nove cilje: C ostane na (3,3), L na (5,5), P pa recimo na (2,2) ali (6,6). Preveri, katere od teh so pri tebi proste?} \\
            }
        }
    }
     \\ \\

    \theutterance \stepcounter{utterance}  
    & & & \multicolumn{2}{p{0.3\linewidth}}{
        \cellcolor[rgb]{0.9,0.9,0.9}{
            \makecell[{{p{\linewidth}}}]{
                \texttt{\tiny{[GM$|$GM]}}
                \texttt{Ponovno preverim ploščo. L je trenutno na (7,3), cilj je (4,4). Na (4,4) je res '├', kar ni prazno. Ciljne lokacije smo določili kot (3,3) za C, (4,4) za L in (5,5) za P. Ampak (4,4) ni prazno. Počakaj, ponovno preverim svojo trenutno ploščo.} \\
\\ 
\texttt{Vrstica 4: `║◌◌◌├─┘P║ 4` {-} stolpci: 1:◌, 2:◌, 3:◌, 4:├, 5:─, 6:┘, 7:P. Torej (4,4) je res zasedeno z '├'. To pomeni, da (4,4) ni dosegljivo. } \\
\\ 
\texttt{Drugi igralec je rekel, da so pri njem te lokacije prazne, ampak pri meni (4,4) ni prazno. Moram predlagati nove ciljne lokacije.} \\
\\ 
\texttt{RECI: Oprosti, na moji plošči (4,4) ni prazno {-} tam je '├'. Prav tako (5,5) preverim... Vrstica 5: `║◌┌─┘◌◌◌║ 5` {-} (5,5) je ◌ (prazno). (3,3) je zdaj C, kar je v redu. Predlagam nove cilje: C ostane na (3,3), L na (5,5), P pa recimo na (2,2) ali (6,6). Preveri, katere od teh so pri tebi proste?} \\
            }
        }
    }
    & & \\ \\

    \theutterance \stepcounter{utterance}  
    & & & \multicolumn{2}{p{0.3\linewidth}}{
        \cellcolor[rgb]{0.9,0.9,0.9}{
            \makecell[{{p{\linewidth}}}]{
                \texttt{\tiny{[GM$|$GM]}}
                \texttt{Invalid format: head is not empty: 'Ponovno preverim ploščo. L je trenutno na (7,3), cilj je (4,4). Na (4,4) je res '├', kar ni prazno. Ciljne lokacije smo določili kot (3,3) za C, (4,4) za L in (5,5) za P. Ampak (4,4) ni prazno. Počakaj, ponovno preverim svojo trenutno ploščo.} \\
\\ 
\texttt{Vrstica 4: ║◌◌◌├─┘P║ 4 {-} stolpci: 1:◌, 2:◌, 3:◌, 4:├, 5:─, 6:┘, 7:P. Torej (4,4) je res zasedeno z '├'. To pomeni, da (4,4) ni dosegljivo. } \\
\\ 
\texttt{Drugi igralec je rekel, da so pri njem te lokacije prazne, ampak pri meni (4,4) ni prazno. Moram predlagati nove ciljne lokacije.} \\
\\ 
\texttt{'} \\
            }
        }
    }
    & & \\ \\

    \theutterance \stepcounter{utterance}  
    & & \multicolumn{4}{p{0.6\linewidth}}{
        \cellcolor[rgb]{0.9,0.9,0.9}{
            \makecell[{{p{\linewidth}}}]{
                \texttt{\tiny{[GM$\rangle$P2]}}
                \texttt{Kazen: Vaše sporočilo ne sme vsebovati ničesar pred ukazom!} \\
\\ 
\texttt{Prepričajte se, da vaš odgovor vsebuje samo `RECI: <SPOROČILO>` ali `PREMAKNI SE: <PREDMET>, (<X>, <Y>)` in nič drugega!} \\
\\ 
\texttt{Skupaj ste zbrali 5 od 8 kazni.} \\
\texttt{Prosim, poskusite znova!} \\
            }
        }
    }
    & \\ \\

    \theutterance \stepcounter{utterance}  
    & & & \multicolumn{4}{p{0.6\linewidth}}{
        \cellcolor[rgb]{0.9,0.9,0.9}{
            \makecell[{{p{\linewidth}}}]{
                \texttt{\tiny{[GM$\langle$P2]}}
                \texttt{RECI: Oprosti, na moji plošči (4,4) ni prazno {-} tam je ovira. Predlagam nove cilje: C ostane na (3,3), L na (5,5), P pa na (2,2) ali (6,6). Preveri, katere so pri tebi proste?} \\
            }
        }
    }
     \\ \\

    \theutterance \stepcounter{utterance}  
    & & & \multicolumn{2}{p{0.3\linewidth}}{
        \cellcolor[rgb]{0.9,0.9,0.9}{
            \makecell[{{p{\linewidth}}}]{
                \texttt{\tiny{[GM$|$GM]}}
                \texttt{RECI: Oprosti, na moji plošči (4,4) ni prazno {-} tam je ovira. Predlagam nove cilje: C ostane na (3,3), L na (5,5), P pa na (2,2) ali (6,6). Preveri, katere so pri tebi proste?} \\
            }
        }
    }
    & & \\ \\

    \theutterance \stepcounter{utterance}  
    & & \multicolumn{4}{p{0.6\linewidth}}{
        \cellcolor[rgb]{0.9,0.9,0.9}{
            \makecell[{{p{\linewidth}}}]{
                \texttt{\tiny{[P1$\langle$GM]}}
                \texttt{Uspešno premaknjen 'C' v (1,1). Vaša posodobljena mreža je videti takole:} \\
\texttt{```} \\
\texttt{ 1234567} \\
\texttt{╔═╤═╤═══╗} \\
\texttt{║C│◌│◌◌◌║ 1} \\
\texttt{╟─┴─┼─┬─╢ 2} \\
\texttt{║◌◌◌│◌│P║ 3} \\
\texttt{║◌◌◌├─┘◌║ 4} \\
\texttt{║◌┌─┘◌◌◌║ 5} \\
\texttt{╟─┤◌◌◌◌◌║ 6} \\
\texttt{║◌│◌L◌◌◌║ 7} \\
\texttt{╚═╧═════╝} \\
\texttt{```} \\
\texttt{Trenutno igrate krog 3 od največ 12.} \\
\texttt{Skupaj ste zbrali 5 od 8 kazni.} \\
\texttt{Drugi igralec je poslal to sporočilo:} \\
\\ 
\\ 
\texttt{"Oprosti, na moji plošči (4,4) ni prazno {-} tam je ovira. Predlagam nove cilje: C ostane na (3,3), L na (5,5), P pa na (2,2) ali (6,6). Preveri, katere so pri tebi proste?"} \\
\\ 
\\ 
\texttt{Kakšen je vaš naslednji ukaz?} \\
            }
        }
    }
    & \\ \\

    \theutterance \stepcounter{utterance}  
    & \multicolumn{4}{p{0.6\linewidth}}{
        \cellcolor[rgb]{0.9,0.9,0.9}{
            \makecell[{{p{\linewidth}}}]{
                \texttt{\tiny{[P1$\rangle$GM]}}
                \texttt{RECI: Pri meni so (3,3), (5,5), (2,2) in (6,6) vse prazne. Predlagam, da izbereva (3,3) za C, (5,5) za L in (6,6) za P. Preveri, če so te tri lokacije na tvoji plošči prazne, in potem začniva premikati!} \\
            }
        }
    }
    & & \\ \\

    \theutterance \stepcounter{utterance}  
    & & & \multicolumn{2}{p{0.3\linewidth}}{
        \cellcolor[rgb]{0.9,0.9,0.9}{
            \makecell[{{p{\linewidth}}}]{
                \texttt{\tiny{[GM$|$GM]}}
                \texttt{RECI: Pri meni so (3,3), (5,5), (2,2) in (6,6) vse prazne. Predlagam, da izbereva (3,3) za C, (5,5) za L in (6,6) za P. Preveri, če so te tri lokacije na tvoji plošči prazne, in potem začniva premikati!} \\
            }
        }
    }
    & & \\ \\

    \theutterance \stepcounter{utterance}  
    & & \multicolumn{4}{p{0.6\linewidth}}{
        \cellcolor[rgb]{0.9,0.9,0.9}{
            \makecell[{{p{\linewidth}}}]{
                \texttt{\tiny{[GM$\rangle$P2]}}
                \texttt{Vaše sporočilo je bilo posredovano drugemu igralcu.} \\
\\ 
\\ 
\texttt{Trenutno igrate krog 4 od največ 12.} \\
\texttt{Skupaj ste zbrali 5 od 8 kazni.} \\
\texttt{Drugi igralec je poslal to sporočilo:} \\
\\ 
\\ 
\texttt{"Pri meni so (3,3), (5,5), (2,2) in (6,6) vse prazne. Predlagam, da izbereva (3,3) za C, (5,5) za L in (6,6) za P. Preveri, če so te tri lokacije na tvoji plošči prazne, in potem začniva premikati!"} \\
\\ 
\\ 
\texttt{Kakšen je vaš naslednji ukaz?} \\
            }
        }
    }
    & \\ \\

    \theutterance \stepcounter{utterance}  
    & & & \multicolumn{4}{p{0.6\linewidth}}{
        \cellcolor[rgb]{0.9,0.9,0.9}{
            \makecell[{{p{\linewidth}}}]{
                \texttt{\tiny{[GM$\langle$P2]}}
                \texttt{Preverim svojo ploščo:} \\
\texttt{{-} (3,3): zdaj je C, v redu} \\
\texttt{{-} (5,5): vrstica 5, stolpec 5 je ◌ (prazno) ✓} \\
\texttt{{-} (6,6): vrstica 6, stolpec 6 je ◌ (prazno) ✓} \\
\\ 
\texttt{Obe lokaciji sta prazni. C je že na mestu. Zdaj premaknem L na (5,5).} \\
\\ 
\texttt{PREMAKNI SE: L, (5,5)} \\
            }
        }
    }
     \\ \\

    \theutterance \stepcounter{utterance}  
    & & & \multicolumn{2}{p{0.3\linewidth}}{
        \cellcolor[rgb]{0.9,0.9,0.9}{
            \makecell[{{p{\linewidth}}}]{
                \texttt{\tiny{[GM$|$GM]}}
                \texttt{Preverim svojo ploščo:} \\
\texttt{{-} (3,3): zdaj je C, v redu} \\
\texttt{{-} (5,5): vrstica 5, stolpec 5 je ◌ (prazno) ✓} \\
\texttt{{-} (6,6): vrstica 6, stolpec 6 je ◌ (prazno) ✓} \\
\\ 
\texttt{Obe lokaciji sta prazni. C je že na mestu. Zdaj premaknem L na (5,5).} \\
\\ 
\texttt{PREMAKNI SE: L, (5,5)} \\
            }
        }
    }
    & & \\ \\

    \theutterance \stepcounter{utterance}  
    & & & \multicolumn{2}{p{0.3\linewidth}}{
        \cellcolor[rgb]{0.9,0.9,0.9}{
            \makecell[{{p{\linewidth}}}]{
                \texttt{\tiny{[GM$|$GM]}}
                \texttt{Invalid format: head is not empty: 'Preverim svojo ploščo:} \\
\texttt{{-} (3,3): zdaj je C, v redu} \\
\texttt{{-} (5,5): vrstica 5, stolpec 5 je ◌ (prazno) ✓} \\
\texttt{{-} (6,6): vrstica 6, stolpec 6 je ◌ (prazno) ✓} \\
\\ 
\texttt{Obe lokaciji sta prazni. C je že na mestu. Zdaj premaknem L na (5,5).} \\
\\ 
\texttt{'} \\
            }
        }
    }
    & & \\ \\

    \theutterance \stepcounter{utterance}  
    & & \multicolumn{4}{p{0.6\linewidth}}{
        \cellcolor[rgb]{0.9,0.9,0.9}{
            \makecell[{{p{\linewidth}}}]{
                \texttt{\tiny{[GM$\rangle$P2]}}
                \texttt{Kazen: Vaše sporočilo ne sme vsebovati ničesar pred ukazom!} \\
\\ 
\texttt{Prepričajte se, da vaš odgovor vsebuje samo `RECI: <SPOROČILO>` ali `PREMAKNI SE: <PREDMET>, (<X>, <Y>)` in nič drugega!} \\
\\ 
\texttt{Skupaj ste zbrali 6 od 8 kazni.} \\
\texttt{Prosim, poskusite znova!} \\
            }
        }
    }
    & \\ \\

    \theutterance \stepcounter{utterance}  
    & & & \multicolumn{4}{p{0.6\linewidth}}{
        \cellcolor[rgb]{0.9,0.9,0.9}{
            \makecell[{{p{\linewidth}}}]{
                \texttt{\tiny{[GM$\langle$P2]}}
                \texttt{PREMAKNI SE: L, (5,5)} \\
            }
        }
    }
     \\ \\

    \theutterance \stepcounter{utterance}  
    & & & \multicolumn{2}{p{0.3\linewidth}}{
        \cellcolor[rgb]{0.9,0.9,0.9}{
            \makecell[{{p{\linewidth}}}]{
                \texttt{\tiny{[GM$|$GM]}}
                \texttt{PREMAKNI SE: L, (5,5)} \\
            }
        }
    }
    & & \\ \\

    \theutterance \stepcounter{utterance}  
    & & & \multicolumn{2}{p{0.3\linewidth}}{
        \cellcolor[rgb]{0.9,0.9,0.9}{
            \makecell[{{p{\linewidth}}}]{
                \texttt{\tiny{[GM$|$GM]}}
                \texttt{L moved to (5, 5)} \\
            }
        }
    }
    & & \\ \\

    \theutterance \stepcounter{utterance}  
    & & \multicolumn{4}{p{0.6\linewidth}}{
        \cellcolor[rgb]{0.9,0.9,0.9}{
            \makecell[{{p{\linewidth}}}]{
                \texttt{\tiny{[P1$\langle$GM]}}
                \texttt{Vaše sporočilo je bilo posredovano drugemu igralcu.} \\
\\ 
\\ 
\texttt{Trenutno igrate krog 4 od največ 12.} \\
\texttt{Skupaj ste zbrali 6 od 8 kazni.} \\
\texttt{Drugi igralec je premaknil predmet na svoji mreži.} \\
\\ 
\\ 
\texttt{Kakšen je vaš naslednji ukaz?} \\
            }
        }
    }
    & \\ \\

    \theutterance \stepcounter{utterance}  
    & \multicolumn{4}{p{0.6\linewidth}}{
        \cellcolor[rgb]{0.9,0.9,0.9}{
            \makecell[{{p{\linewidth}}}]{
                \texttt{\tiny{[P1$\rangle$GM]}}
                \texttt{PREMAKNI SE: P, (6,3)} \\
            }
        }
    }
    & & \\ \\

    \theutterance \stepcounter{utterance}  
    & & & \multicolumn{2}{p{0.3\linewidth}}{
        \cellcolor[rgb]{0.9,0.9,0.9}{
            \makecell[{{p{\linewidth}}}]{
                \texttt{\tiny{[GM$|$GM]}}
                \texttt{PREMAKNI SE: P, (6,3)} \\
            }
        }
    }
    & & \\ \\

    \theutterance \stepcounter{utterance}  
    & & & \multicolumn{2}{p{0.3\linewidth}}{
        \cellcolor[rgb]{0.9,0.9,0.9}{
            \makecell[{{p{\linewidth}}}]{
                \texttt{\tiny{[GM$|$GM]}}
                \texttt{Kazen: (6,3) ni prazen, vendar vsebuje '│'.} \\
\\ 
\texttt{Skupaj ste zbrali 7 od 8 kazni.} \\
\texttt{Prosim, poskusite znova!} \\
            }
        }
    }
    & & \\ \\

    \theutterance \stepcounter{utterance}  
    & & \multicolumn{4}{p{0.6\linewidth}}{
        \cellcolor[rgb]{0.9,0.9,0.9}{
            \makecell[{{p{\linewidth}}}]{
                \texttt{\tiny{[P1$\langle$GM]}}
                \texttt{Kazen: (6,3) ni prazen, vendar vsebuje '│'.} \\
\\ 
\texttt{Skupaj ste zbrali 7 od 8 kazni.} \\
\texttt{Prosim, poskusite znova!} \\
            }
        }
    }
    & \\ \\

    \theutterance \stepcounter{utterance}  
    & \multicolumn{4}{p{0.6\linewidth}}{
        \cellcolor[rgb]{0.9,0.9,0.9}{
            \makecell[{{p{\linewidth}}}]{
                \texttt{\tiny{[P1$\rangle$GM]}}
                \texttt{PREMAKNI SE: P, (7,2)} \\
            }
        }
    }
    & & \\ \\

    \theutterance \stepcounter{utterance}  
    & & & \multicolumn{2}{p{0.3\linewidth}}{
        \cellcolor[rgb]{0.9,0.9,0.9}{
            \makecell[{{p{\linewidth}}}]{
                \texttt{\tiny{[GM$|$GM]}}
                \texttt{PREMAKNI SE: P, (7,2)} \\
            }
        }
    }
    & & \\ \\

    \theutterance \stepcounter{utterance}  
    & & & \multicolumn{2}{p{0.3\linewidth}}{
        \cellcolor[rgb]{0.9,0.9,0.9}{
            \makecell[{{p{\linewidth}}}]{
                \texttt{\tiny{[GM$|$GM]}}
                \texttt{Kazen: (7,2) ni prazen, vendar vsebuje '─'.} \\
\\ 
\texttt{Skupaj ste zbrali 8 od 8 kazni.} \\
\texttt{Prosim, poskusite znova!} \\
            }
        }
    }
    & & \\ \\

    \theutterance \stepcounter{utterance}  
    & & \multicolumn{4}{p{0.6\linewidth}}{
        \cellcolor[rgb]{0.9,0.9,0.9}{
            \makecell[{{p{\linewidth}}}]{
                \texttt{\tiny{[P1$\langle$GM]}}
                \texttt{Kazen: (7,2) ni prazen, vendar vsebuje '─'.} \\
\\ 
\texttt{Skupaj ste zbrali 8 od 8 kazni.} \\
\texttt{Prosim, poskusite znova!} \\
            }
        }
    }
    & \\ \\

    \theutterance \stepcounter{utterance}  
    & \multicolumn{4}{p{0.6\linewidth}}{
        \cellcolor[rgb]{0.9,0.9,0.9}{
            \makecell[{{p{\linewidth}}}]{
                \texttt{\tiny{[P1$\rangle$GM]}}
                \texttt{PREMAKNI SE: P, (6,2)} \\
            }
        }
    }
    & & \\ \\

    \theutterance \stepcounter{utterance}  
    & & & \multicolumn{2}{p{0.3\linewidth}}{
        \cellcolor[rgb]{0.9,0.9,0.9}{
            \makecell[{{p{\linewidth}}}]{
                \texttt{\tiny{[GM$|$GM]}}
                \texttt{PREMAKNI SE: P, (6,2)} \\
            }
        }
    }
    & & \\ \\

    \theutterance \stepcounter{utterance}  
    & & & \multicolumn{2}{p{0.3\linewidth}}{
        \cellcolor[rgb]{0.9,0.9,0.9}{
            \makecell[{{p{\linewidth}}}]{
                \texttt{\tiny{[GM$|$GM]}}
                \texttt{Kazen: (6,2) ni prazen, vendar vsebuje '┬'.} \\
\\ 
\texttt{Skupaj ste zbrali 9 od 8 kazni.} \\
\texttt{Prosim, poskusite znova!} \\
            }
        }
    }
    & & \\ \\

    \theutterance \stepcounter{utterance}  
    & & \multicolumn{4}{p{0.6\linewidth}}{
        \cellcolor[rgb]{0.9,0.9,0.9}{
            \makecell[{{p{\linewidth}}}]{
                \texttt{\tiny{[P1$\langle$GM]}}
                \texttt{Kazen: (6,2) ni prazen, vendar vsebuje '┬'.} \\
\\ 
\texttt{Skupaj ste zbrali 9 od 8 kazni.} \\
\texttt{Prosim, poskusite znova!} \\
            }
        }
    }
    & \\ \\

    \theutterance \stepcounter{utterance}  
    & \multicolumn{4}{p{0.6\linewidth}}{
        \cellcolor[rgb]{0.9,0.9,0.9}{
            \makecell[{{p{\linewidth}}}]{
                \texttt{\tiny{[P1$\rangle$GM]}}
                \texttt{PREMAKNI SE: P, (5,3)} \\
            }
        }
    }
    & & \\ \\

    \theutterance \stepcounter{utterance}  
    & & & \multicolumn{2}{p{0.3\linewidth}}{
        \cellcolor[rgb]{0.9,0.9,0.9}{
            \makecell[{{p{\linewidth}}}]{
                \texttt{\tiny{[GM$|$GM]}}
                \texttt{PREMAKNI SE: P, (5,3)} \\
            }
        }
    }
    & & \\ \\

    \theutterance \stepcounter{utterance}  
    & & & \multicolumn{2}{p{0.3\linewidth}}{
        \cellcolor[rgb]{0.9,0.9,0.9}{
            \makecell[{{p{\linewidth}}}]{
                \texttt{\tiny{[GM$|$GM]}}
                \texttt{P moved to (5, 3)} \\
            }
        }
    }
    & & \\ \\

    \theutterance \stepcounter{utterance}  
    & & \multicolumn{4}{p{0.6\linewidth}}{
        \cellcolor[rgb]{0.9,0.9,0.9}{
            \makecell[{{p{\linewidth}}}]{
                \texttt{\tiny{[GM$\rangle$P2]}}
                \texttt{Uspešno premaknjen 'L' v (5,5). Vaša posodobljena mreža je videti takole:} \\
\texttt{```} \\
\texttt{ 1234567} \\
\texttt{╔═╤═╤═══╗} \\
\texttt{║◌│◌│◌◌◌║ 1} \\
\texttt{╟─┴─┼─┬─╢ 2} \\
\texttt{║◌◌C│◌│◌║ 3} \\
\texttt{║◌◌◌├─┘P║ 4} \\
\texttt{║◌┌─┘L◌◌║ 5} \\
\texttt{╟─┤◌◌◌◌◌║ 6} \\
\texttt{║◌│◌◌◌◌◌║ 7} \\
\texttt{╚═╧═════╝} \\
\texttt{```} \\
\texttt{Trenutno igrate krog 5 od največ 12.} \\
\texttt{Skupaj ste zbrali 9 od 8 kazni.} \\
\texttt{Drugi igralec je premaknil predmet na svoji mreži.} \\
\\ 
\\ 
\texttt{Kakšen je vaš naslednji ukaz?} \\
            }
        }
    }
    & \\ \\

    \theutterance \stepcounter{utterance}  
    & & & \multicolumn{4}{p{0.6\linewidth}}{
        \cellcolor[rgb]{0.9,0.9,0.9}{
            \makecell[{{p{\linewidth}}}]{
                \texttt{\tiny{[GM$\langle$P2]}}
                \texttt{Kazni so presegle 8, oba sva izgubila igro. Ampak vidim, da igra še poteka, morda sistem še ni končal. Poskusim premakniti P na (6,6).} \\
\\ 
\texttt{PREMAKNI SE: P, (6,6)} \\
            }
        }
    }
     \\ \\

    \theutterance \stepcounter{utterance}  
    & & & \multicolumn{2}{p{0.3\linewidth}}{
        \cellcolor[rgb]{0.9,0.9,0.9}{
            \makecell[{{p{\linewidth}}}]{
                \texttt{\tiny{[GM$|$GM]}}
                \texttt{Kazni so presegle 8, oba sva izgubila igro. Ampak vidim, da igra še poteka, morda sistem še ni končal. Poskusim premakniti P na (6,6).} \\
\\ 
\texttt{PREMAKNI SE: P, (6,6)} \\
            }
        }
    }
    & & \\ \\

    \theutterance \stepcounter{utterance}  
    & & & \multicolumn{2}{p{0.3\linewidth}}{
        \cellcolor[rgb]{0.9,0.9,0.9}{
            \makecell[{{p{\linewidth}}}]{
                \texttt{\tiny{[GM$|$GM]}}
                \texttt{Invalid format: head is not empty: 'Kazni so presegle 8, oba sva izgubila igro. Ampak vidim, da igra še poteka, morda sistem še ni končal. Poskusim premakniti P na (6,6).} \\
\\ 
\texttt{'} \\
            }
        }
    }
    & & \\ \\

    \theutterance \stepcounter{utterance}  
    & & & \multicolumn{2}{p{0.3\linewidth}}{
        \cellcolor[rgb]{0.9,0.9,0.9}{
            \makecell[{{p{\linewidth}}}]{
                \texttt{\tiny{[GM$|$GM]}}
                \texttt{Maximum number of penalties exceeded} \\
            }
        }
    }
    & & \\ \\

    \theutterance \stepcounter{utterance}  
    & & & \multicolumn{2}{p{0.3\linewidth}}{
        \cellcolor[rgb]{0.9,0.9,0.9}{
            \makecell[{{p{\linewidth}}}]{
                \texttt{\tiny{[GM$|$GM]}}
                \texttt{Initial states:} \\
\texttt{```} \\
\texttt{Player 1:} \\
\texttt{ 1234567} \\
\texttt{╔═╤═╤═══╗} \\
\texttt{║◌│◌│C◌◌║ 1} \\
\texttt{╟─┴─┼─┬─╢ 2} \\
\texttt{║◌◌◌│◌│P║ 3} \\
\texttt{║◌◌◌├─┘◌║ 4} \\
\texttt{║◌┌─┘◌◌◌║ 5} \\
\texttt{╟─┤◌◌◌◌◌║ 6} \\
\texttt{║◌│◌L◌◌◌║ 7} \\
\texttt{╚═╧═════╝} \\
\\ 
\\ 
\texttt{Player 2:} \\
\texttt{ 1234567} \\
\texttt{╔═╤═╤═══╗} \\
\texttt{║◌│◌│◌◌◌║ 1} \\
\texttt{╟─┴─┼─┬─╢ 2} \\
\texttt{║◌◌◌│◌│L║ 3} \\
\texttt{║C◌◌├─┘P║ 4} \\
\texttt{║◌┌─┘◌◌◌║ 5} \\
\texttt{╟─┤◌◌◌◌◌║ 6} \\
\texttt{║◌│◌◌◌◌◌║ 7} \\
\texttt{╚═╧═════╝} \\
\\ 
\texttt{```} \\
            }
        }
    }
    & & \\ \\

    \theutterance \stepcounter{utterance}  
    & & & \multicolumn{2}{p{0.3\linewidth}}{
        \cellcolor[rgb]{0.9,0.9,0.9}{
            \makecell[{{p{\linewidth}}}]{
                \texttt{\tiny{[GM$|$GM]}}
                \texttt{End states:} \\
\texttt{```} \\
\texttt{Player 1:} \\
\texttt{ 1234567} \\
\texttt{╔═╤═╤═══╗} \\
\texttt{║C│◌│◌◌◌║ 1} \\
\texttt{╟─┴─┼─┬─╢ 2} \\
\texttt{║◌◌◌│P│◌║ 3} \\
\texttt{║◌◌◌├─┘◌║ 4} \\
\texttt{║◌┌─┘◌◌◌║ 5} \\
\texttt{╟─┤◌◌◌◌◌║ 6} \\
\texttt{║◌│◌L◌◌◌║ 7} \\
\texttt{╚═╧═════╝} \\
\\ 
\\ 
\texttt{Player 2:} \\
\texttt{ 1234567} \\
\texttt{╔═╤═╤═══╗} \\
\texttt{║◌│◌│◌◌◌║ 1} \\
\texttt{╟─┴─┼─┬─╢ 2} \\
\texttt{║◌◌C│◌│◌║ 3} \\
\texttt{║◌◌◌├─┘P║ 4} \\
\texttt{║◌┌─┘L◌◌║ 5} \\
\texttt{╟─┤◌◌◌◌◌║ 6} \\
\texttt{║◌│◌◌◌◌◌║ 7} \\
\texttt{╚═╧═════╝} \\
\\ 
\texttt{```} \\
            }
        }
    }
    & & \\ \\

    \theutterance \stepcounter{utterance}  
    & & & \multicolumn{2}{p{0.3\linewidth}}{
        \cellcolor[rgb]{0.9,0.9,0.9}{
            \makecell[{{p{\linewidth}}}]{
                \texttt{\tiny{[GM$|$GM]}}
                \texttt{* success: False} \\
\texttt{* lose: False} \\
\texttt{* aborted: True} \\
\texttt{{-}{-}{-}{-}{-}{-}{-}} \\
\texttt{* Shifts: 2.00} \\
\texttt{* Max Shifts: 4.00} \\
\texttt{* Min Shifts: 2.00} \\
\texttt{* End Distance Sum: 7.30} \\
\texttt{* Init Distance Sum: 11.00} \\
\texttt{* Expected Distance Sum: 12.57} \\
\texttt{* Penalties: 10.00} \\
\texttt{* Max Penalties: 8.00} \\
\texttt{* Rounds: 4.00} \\
\texttt{* Max Rounds: 12.00} \\
\texttt{* Object Count: 3.00} \\
            }
        }
    }
    & & \\ \\

\end{supertabular}
}

\end{document}
